\clearpage
\begin{explanationbox}
\textbf{\textcolor{CExplanation}{一、为什么会想到这一结论？}}
\begin{center}\begin{tikzpicture}[font=\small,>=stealth,line join=round]
\draw[gray!40,rounded corners=3pt] (0,0) rectangle (14.2,8.3);\node at (7.1,7.700000000000001) {两个输入相隔4，所以每个位置都能重复};
\begin{scope}[shift={(1.882608695652174,3.7409090909090907)},x=0.5217391304347826cm,y=1.3939393939393938cm]\draw[->,gray] (-1.5,0)--(10,0) node[right]{$x$};\draw[->,gray] (0,-1.5)--(0,1.8) node[above]{$y$};\draw[blue,thick,solid,domain=-1.2:9.6,samples=180] plot(\x,{sin(deg(pi*\x/2))});\fill[orange!85!black] (1,1) circle(2.2pt);\fill[orange!85!black] (5,1) circle(2.2pt);\fill[orange!85!black] (9,1) circle(2.2pt);\foreach \v in {1,5,9}{\draw[gray,dashed] (\v,0)--(\v,1);\node[below=6pt] at (\v,0){$\v$};}\end{scope}\node[anchor=west,align=left,text width=5.0cm] at (8.4,4.0) {示例 $f=\sin\frac{\pi x}2$\\[9pt]$a=2$\\[9pt]$f(1)=f(5)=f(9)=1$\\[9pt]$f(x+2)=f(x-2)$\\[9pt]周期 $T=4$};
\node[text width=12.6cm,align=center] at (7.1,.65) {曲线示范整段重复；一般证明依赖每一个输入都满足等值关系。};
\end{tikzpicture}\end{center}



我们熟悉“函数每隔 $T$ 就重复一次”的写法：$f(x+T)=f(x)$。可是考试不总把它写得这么整齐，有时给出的是
\[
f(x+a)=f(x-a).
\]
这就像问：\textbf{两个位置的高度总是一样，说明图像按多长的间隔在重复？}
先不背公式，先在数轴上找出这两个位置。

\medskip

\end{explanationbox}
\clearpage
\begin{explanationbox}
\textbf{\textcolor{CExplanation}{二、初中生也能看懂：用数字把距离算出来}}
\begin{center}\begin{tikzpicture}[font=\small,>=stealth,line join=round]
\draw[gray!40,rounded corners=3pt] (0,0) rectangle (14.2,8.3);\node at (7.1,7.700000000000001) {算的是两个输入之差，不是中点到单侧的距离};
\draw[->,gray] (.8,3.2)--(13.1,3.2) node[right]{$x$};\foreach \v in {1,3,5,7,9}{\pgfmathsetmacro{\X}{1.8+\v}\draw[gray] (\X,3.2)--(\X,3.1);\node[below=7pt] at (\X,3.1){$\v$};}\foreach \v in {1,5,9}{\fill[blue] ({1.8+\v},3.2) circle(2.4pt);}\draw[<->,orange!85!black,thick] (2.8,4.7)--(6.8,4.7) node[midway,above]{$5-1=4$};\draw[<->,violet,thick] (6.8,6.1)--(10.8,6.1) node[midway,above]{$9-5=4$};\node at (7.1,1.8){$x=3$ 时，输入为 $x-2=1$ 与 $x+2=5$；中点到两边各2。};
\node[text width=12.6cm,align=center] at (7.1,.65) {相邻标点相隔4；仅凭这几个点，不能证明整个函数的周期。};
\end{tikzpicture}\end{center}



设 $a=2$，条件就成了 $f(x+2)=f(x-2)$。
\begin{itemize}[leftmargin=*,itemsep=0.3em]
\item 把 $x=3$ 代入：$f(5)=f(1)$。两个位置 $1$ 和 $5$ 相隔 $4$。
\item 再把 $x=7$ 代入：$f(9)=f(5)$。两个位置 $5$ 和 $9$ 也相隔 $4$。
\item 所以 $f(1)=f(5)=f(9)$。不是只在这三个位置巧合；\textbf{由于题目对任意 $x$ 成立，每个位置都可以这样往右平移 $4$。}
\end{itemize}

\[
\underbrace{(x+a)}_{\text{右边位置}}-\underbrace{(x-a)}_{\text{左边位置}}=2a.
\]
当 $a>0$ 时，从 $x-a$ 到 $x+a$ 要走 $2a$；当 $a<0$ 时两边左右互换，实际距离为 $2|a|$。

\medskip

\end{explanationbox}
\clearpage
\begin{explanationbox}
\textbf{\textcolor{CExplanation}{三、数轴图解：为什么不是 $a$？}}
\begin{center}\begin{tikzpicture}[font=\small,>=stealth,line join=round]
\draw[gray!40,rounded corners=3pt] (0,0) rectangle (14.2,8.3);\node at (7.1,7.700000000000001) {算的是两个输入之差，不是中点到单侧的距离};
\draw[->,gray] (.8,3.2)--(13.1,3.2) node[right]{$x$};\foreach \v in {1,3,5,7,9}{\pgfmathsetmacro{\X}{1.8+\v}\draw[gray] (\X,3.2)--(\X,3.1);\node[below=7pt] at (\X,3.1){$\v$};}\foreach \v in {1,5,9}{\fill[blue] ({1.8+\v},3.2) circle(2.4pt);}\draw[<->,orange!85!black,thick] (2.8,4.7)--(6.8,4.7) node[midway,above]{$5-1=4$};\draw[<->,violet,thick] (6.8,6.1)--(10.8,6.1) node[midway,above]{$9-5=4$};\node at (7.1,1.8){$x=3$ 时，输入为 $x-2=1$ 与 $x+2=5$；中点到两边各2。};
\node[text width=12.6cm,align=center] at (7.1,.65) {相邻标点相隔4；仅凭这几个点，不能证明整个函数的周期。};
\end{tikzpicture}\end{center}



下面的 $x=3$ 只是用来帮助看图；\textbf{使整个函数有周期的是“任意 $x$ 都成立”}，不是只知道某三个点一样高。


\textbf{读图说明：} 当 $x=3$ 时，比较的是 $x-2=1$ 和 $x+2=5$；换成别的 $x$，仍然是比较两处相隔 $4$ 的位置。

\medskip

\end{explanationbox}
\clearpage
\begin{explanationbox}
\textbf{\textcolor{CExplanation}{四、把它画成真正的函数图像}}
\begin{center}\begin{tikzpicture}[font=\small,>=stealth,line join=round]
\draw[gray!40,rounded corners=3pt] (0,0) rectangle (14.2,8.3);\node at (7.1,7.700000000000001) {两个输入相隔4，所以每个位置都能重复};
\begin{scope}[shift={(1.882608695652174,3.7409090909090907)},x=0.5217391304347826cm,y=1.3939393939393938cm]\draw[->,gray] (-1.5,0)--(10,0) node[right]{$x$};\draw[->,gray] (0,-1.5)--(0,1.8) node[above]{$y$};\draw[blue,thick,solid,domain=-1.2:9.6,samples=180] plot(\x,{sin(deg(pi*\x/2))});\fill[orange!85!black] (1,1) circle(2.2pt);\fill[orange!85!black] (5,1) circle(2.2pt);\fill[orange!85!black] (9,1) circle(2.2pt);\foreach \v in {1,5,9}{\draw[gray,dashed] (\v,0)--(\v,1);\node[below=6pt] at (\v,0){$\v$};}\end{scope}\node[anchor=west,align=left,text width=5.0cm] at (8.4,4.0) {示例 $f=\sin\frac{\pi x}2$\\[9pt]$a=2$\\[9pt]$f(1)=f(5)=f(9)=1$\\[9pt]$f(x+2)=f(x-2)$\\[9pt]周期 $T=4$};
\node[text width=12.6cm,align=center] at (7.1,.65) {曲线示范整段重复；一般证明依赖每一个输入都满足等值关系。};
\end{tikzpicture}\end{center}



例如 $f(x)=\sin\!\left(\dfrac{\pi x}{2}\right)$。这条曲线在 $x=1,5,9$ 都取值 $1$，\textbf{相邻两段图像的形状完全相同}，不只是几个点一样高。


\textbf{注意：} 这是为了帮助想象而选的一个具体函数；证明任意满足条件的函数有周期，不需要知道它的解析式。

\medskip

\end{explanationbox}
\clearpage
\begin{explanationbox}
\textbf{\textcolor{CExplanation}{五、真正关键的一步：把“左右两边”改成“从某点往右”}}
\begin{center}\begin{tikzpicture}[font=\small,>=stealth,line join=round]
\draw[gray!40,rounded corners=3pt] (0,0) rectangle (14.2,8.3);\node at (7.1,7.700000000000001) {算的是两个输入之差，不是中点到单侧的距离};
\draw[->,gray] (.8,3.2)--(13.1,3.2) node[right]{$x$};\foreach \v in {1,3,5,7,9}{\pgfmathsetmacro{\X}{1.8+\v}\draw[gray] (\X,3.2)--(\X,3.1);\node[below=7pt] at (\X,3.1){$\v$};}\foreach \v in {1,5,9}{\fill[blue] ({1.8+\v},3.2) circle(2.4pt);}\draw[<->,orange!85!black,thick] (2.8,4.7)--(6.8,4.7) node[midway,above]{$5-1=4$};\draw[<->,violet,thick] (6.8,6.1)--(10.8,6.1) node[midway,above]{$9-5=4$};\node at (7.1,1.8){$x=3$ 时，输入为 $x-2=1$ 与 $x+2=5$；中点到两边各2。};
\node[text width=12.6cm,align=center] at (7.1,.65) {相邻标点相隔4；仅凭这几个点，不能证明整个函数的周期。};
\end{tikzpicture}\end{center}



把 $x-a$ 暂时叫作 $u$，那么 $x+a=u+2a$。于是原式立刻变成
\[
f(u+2a)=f(u).
\]
这才是周期的标准写法。\textbf{周期不是看 $x$ 自己走了多少，而是看两次函数值相等时，两个输入相差多少。}

\medskip

\end{explanationbox}
\clearpage
\begin{explanationbox}
\textbf{\textcolor{CExplanation}{六、眼前一亮：两条对称轴也能推出周期}}
\begin{center}\begin{tikzpicture}[font=\small,>=stealth,line join=round]
\draw[gray!40,rounded corners=3pt] (0,0) rectangle (14.2,8.3);\node at (7.1,7.700000000000001) {两条轴分别配对，两次反射得到水平平移6};
\begin{scope}[shift={(1.9181818181818182,3.7409090909090907)},x=0.5454545454545454cm,y=1.3939393939393938cm]\draw[->,gray] (-1.5,0)--(9.5,0) node[right]{$x$};\draw[->,gray] (0,-1.5)--(0,1.8) node[above]{$y$};\draw[blue,thick,solid,domain=-1:9,samples=180] plot(\x,{cos(deg(pi*(\x-1)/3))});\fill[orange!85!black] (0,0.5) circle(2.2pt);\fill[orange!85!black] (2,0.5) circle(2.2pt);\fill[orange!85!black] (6,0.5) circle(2.2pt);\foreach \a in {1,4}{\draw[violet,dashed] (\a,-1.4)--(\a,1.4);\node[above,violet] at (\a,1.4){$x=\a$};}\end{scope}\node[anchor=west,align=left,text width=5.0cm] at (8.4,4.0) {示例 $f=\cos\frac{\pi(x-1)}3$\\[9pt]轴 $x=1,\ x=4$\\[9pt]$0\to2\to6$\\[9pt]高度都为 $\frac12$\\[9pt]周期 $T=2(4-1)=6$};
\node[text width=12.6cm,align=center] at (7.1,.65) {只讨论完整图像的两条轴；得出的6是一个周期，不一定最小。};
\end{tikzpicture}\end{center}



如果一条完整函数图像同时关于两条不同的竖直直线 $x=p$ 和 $x=q$ 对称：
先关于 $x=p$ 翻过去，再关于 $x=q$ 翻过去，\textbf{效果相当于把整张图沿水平方向挪动 $2(q-p)$}。
由于两次翻折后仍是原图，故
\[
\boxed{T=2|q-p|\text{ 是一个周期。}}
\]
例如，对称轴是 $x=1$ 与 $x=4$，两轴相隔 $3$，可得一个正周期 $6$。\textbf{这是连接“图像对称”与“函数周期”的桥梁。} 注意：这里要求整张图像分别关于这两条直线对称，并非随便挑两个点的函数值相等。

\medskip

\end{explanationbox}
\clearpage
\begin{explanationbox}
\textbf{\textcolor{CExplanation}{七、与已学知识的联系}}
\begin{center}\begin{tikzpicture}[font=\small,>=stealth,line join=round]
\draw[gray!40,rounded corners=3pt] (0,0) rectangle (14.2,8.5);\node at (7.1,7.9) {固定中点的等高是轴对称；比较中点变化才对应此处周期};
\node at (2.4,6.9) {轴对称：$f=(x-2)^2$};\begin{scope}[shift={(1.5499999999999998,3.628571428571429)},x=0.425cm,y=0.5142857142857143cm]\draw[->,gray] (-2,0)--(6,0) node[right]{$x$};\draw[->,gray] (0,-2)--(0,5) node[above]{$y$};\draw[blue,thick,solid,domain=0:4,samples=160] plot(\x,{(\x-2)^2});\fill[orange!85!black] (1,1) circle(2pt);\fill[orange!85!black] (3,1) circle(2pt);\end{scope}\node[align=center,text width=4cm] at (2.4,1.7) {固定中点2：\\$f(1)=f(3)$};\node at (7.050000000000001,6.9) {周期示例：$g=\sin\frac{\pi x}2$};\begin{scope}[shift={(6.2,3.628571428571429)},x=0.425cm,y=0.5142857142857143cm]\draw[->,gray] (-2,0)--(6,0) node[right]{$x$};\draw[->,gray] (0,-2)--(0,5) node[above]{$y$};\draw[violet,thick,solid,domain=-1.8:5.8,samples=160] plot(\x,{sin(deg(pi*\x/2))});\fill[orange!85!black] (-1.5,-0.7071067811865476) circle(2pt);\fill[orange!85!black] (2.5,-0.7071067811865476) circle(2pt);\end{scope}\node[align=center,text width=4cm] at (7.050000000000001,1.7) {中点0.5：\\$g(-1.5)=g(2.5)$};\node at (11.7,6.9) {同一周期函数 $g$};\begin{scope}[shift={(10.85,3.628571428571429)},x=0.425cm,y=0.5142857142857143cm]\draw[->,gray] (-2,0)--(6,0) node[right]{$x$};\draw[->,gray] (0,-2)--(0,5) node[above]{$y$};\draw[violet,thick,solid,domain=-1.8:5.8,samples=160] plot(\x,{sin(deg(pi*\x/2))});\fill[orange!85!black] (0.5,0.7071067811865476) circle(2pt);\fill[orange!85!black] (4.5,0.7071067811865476) circle(2pt);\end{scope}\node[align=center,text width=4cm] at (11.7,1.7) {中点2.5：\\$g(0.5)=g(4.5)$};
\node[text width=12.6cm,align=center] at (7.1,.65) {三图同尺度；右侧两例的比较中点并非对称轴，配对点之间始终相距4。};
\end{tikzpicture}\end{center}


\begin{itemize}[leftmargin=*,itemsep=.35em]
\item $f(x+T)=f(x)$：直接告诉你每走 $T$ 就重复。
\item $f(x+a)=f(x-a)$：看似一左一右，实际两位置相隔 $2|a|$。
\item $f(x+p)=f(x+q)$：更一般的写法，两位置相隔 $|p-q|$。
\item $f(c+t)=f(c-t)$：\textbf{固定中点 $c$} 的两边一样高，表示关于 $x=c$ 对称；与上面的“每一个中点都比较”不是一回事。
\end{itemize}
\end{explanationbox}
